\chapter{Glossary}\label{gloss:chapter}

Terms are listed alphabetically. Each entry gives the v0.1 meaning as used in this monograph and a plain-text pointer to the chapter where the term is treated normatively; where a term is defined in one chapter and stressed in another, both are named. Definitions here are summaries for orientation --- the named chapter's text governs. All defined artifacts and mechanisms are Ring-4 research candidate material.

\begin{description}

\item[Advisory lock.] The profile-bound lock through which cooperating compatibility workers serialize execution: 90-second acquisition bound, 10-millisecond poll interval, 10-millisecond post-release scheduler yield, no retry of a failed worker, no fairness claim. Binds only processes that cooperate. (Chapters: Concurrency and Failure Discipline; Threat Model, lock-based denial.)

\item[Alignment link.] One of the three link kinds; asserts a reviewed structural correspondence between references in different namespaces, weaker than identity and stronger than a hypothesis. Carries its own source and review state. (Chapters: Semantic Identity; Integration Contracts.)

\item[Ambient evidence.] Evidence in its local, context-bound form: a record true of the host, profile, time, and boundary it names, with no cross-context force. All Core records are ambient; resolution is future topos-side work. (Chapter: Integration Contracts, Noetic Topos section.)

\item[Admission state.] First field of the research boundary; records where an artifact stands in the program's admission process. \textsc{open} is the only Ring-4 value. (Chapter: The Research Boundary.)

\item[\texttt{allowed}.] The compatibility protocol's legacy spelling of the authorization field; constant false in v0.1, jointly with \texttt{authorization\_granted}. (Chapter: Runtime Decision Model.)

\item[\texttt{authorization\_granted}.] The decision model's authorization field. Constant false in v0.1; no code path assigns true; claim-bearing intent is denied even when a locally constructed ledger reads \textsc{licensed}. (Chapters: Runtime Decision Model; Evidence, Obligations, and Program.)

\item[Boundary (research boundary).] The jointly validated triple (admission state, product state, public disclosure). Ring 4 requires \textsc{open}, \textsc{not licensed}, disclosure false; contradictory combinations are invalid at construction; every public operation derives and rechecks it. (Chapter: The Research Boundary.)

\item[Bounded canonicalization.] The property that canonical encoding fails closed on structures exceeding declared depth or node limits and on any cycle, guaranteeing termination and bounded resource use. (Chapter: Canonical Evidence Encoding.)

\item[Byte-backed loading.] The registry discipline under which parsed headers, schemas, and records are reconstructed from and rechecked against retained raw ledger and schema bytes, so no constructor path can pair an old raw identity with substituted parsed content. (Chapters: Registry Integrity; Threat Model, registry substitution.)

\item[Byte identity.] The dual identity recorded for tracked files: raw digest of exact bytes plus the CRLF-to-LF comparison digest, with lone-CR, CRLF, and lone-LF counts kept unconflated. Neither digest substitutes for the other. (Chapters: Canonical Evidence Encoding; Threat Model, newline rewriting.)

\item[Capability fields.] Explicit false-valued fields in every compatibility receipt (hermetic execution false; OS-level sandbox false) that state the protocol's limitations inside the artifact itself, so no consumer can honestly over-read it. (Chapters: Worker Protocol; Evidence, Obligations, and Program.)

\item[CAS-Ohm 1.0.] Companion candidate for measurement and instrumentation content: typed electrical quantities, calibration profiles, fail-closed research packets. Mirrors Core's authorization split as completeness versus eligibility. (Chapter: Integration Contracts.)

\item[CAS PCBA 8.] Companion candidate for electronics-engineering content: board model, design-rule checks, manufacturing and test traces. The heaviest consumer of the Core contract surface. (Chapter: Integration Contracts.)

\item[Claim-bearing intent.] The effect class of an action that would assert, promote, or rely on a claim's status. Always denied in v0.1, regardless of structural evidence. Contrast: nonclaim research evidence. (Chapter: Runtime Decision Model.)

\item[Claim record / claim registry.] A registry record registering a claim's exact statement, ring, and boundary; the registry is the byte-backed ledger of such records. At the Gate-4 receipt: eleven records, all Ring 4, all \textsc{not licensed}. (Chapters: Registry Integrity; Evidence, Obligations, and Program.)

\item[Comparison identity.] The declared CRLF-to-LF normalized identity of a file, recorded beside (never instead of) its raw identity. The only normalization the kernel declares. (Chapter: Canonical Evidence Encoding.)

\item[Compatibility profile.] The configuration under which legacy execution is admitted. The only v0.1 profile is \texttt{legacy\_only}: it binds the exact active Python CAS-notation-compiler bytes and produces no v8 semantic translation. (Chapter: Worker Protocol.)

\item[Compatibility receipt.] The record produced by a deterministic-closure execution: source and compiler identities, engine and policy profiles, worker runtime binding, environment inventory, result digest, concurrency record, capability fields, boundary, replay linkage. Historical software evidence, never authority. (Chapters: Worker Protocol; Reference Tables.)

\item[Completeness versus eligibility.] CAS-Ohm's mirror of the authorization split: a packet is complete when structurally whole, eligible only under external authenticated status; completeness never implies eligibility, and eligibility is constant false in v1.0. (Chapter: Integration Contracts.)

\item[Conformance campaign.] The open differential-testing program required before Python/Rust agreement may be asserted: shared hash-bound corpus, digest-level comparison, receipted results, specification-erratum stream. Open obligation two. (Chapter: Evidence, Obligations, and Program.)

\item[Constant-false invariant.] The design stance that \texttt{authorization\_granted} and \texttt{allowed} are false on every path in v0.1; the kernel's terminal defense against receipt laundering. (Chapters: Runtime Decision Model; Threat Model.)

\item[Contract surface.] The eight consumable contracts of Core v0.1 (semantic identity; exact-type discipline; boundary; canonical encoding; byte identity; compatibility receipts and bridge; registry integrity; decision model). Companions consume subsets; none may extend a contract in place. (Chapter: Integration Contracts.)

\item[Decision record.] Output of the runtime decision model: enum-typed action description plus verdict. A structural pass requires fresh replay agreement; claim-bearing intent is denied. (Chapter: Runtime Decision Model.)

\item[\texttt{definition\_sha256}.] Fifth field of a semantic reference: the digest of the exact defining bytes in the identified source, making definitional drift observable at the identity layer. (Chapter: Semantic Identity.)

\item[Dependency direction (strict).] The integration invariant: companions depend on Core contracts; Core depends on no companion semantics. Grounds the reviewability of the kernel in isolation. (Chapter: Integration Contracts.)

\item[Dependency profile.] The hashed description of a worker's effective environment (module inventory, search paths, platform inputs) used to detect divergence between workers. (Chapters: Concurrency and Failure Discipline; Threat Model, divergence.)

\item[Deterministic closure.] The eight-step compatibility protocol: verify source hash; hash compiler bytes before execution; force engine profile and frozen policy; reject \texttt{use} imports; execute in a fresh sanitized worker; bind worker runtime pre-execution; serialize via the advisory lock; record identities and boundary; require fresh replay before structural decisions. (Chapter: Worker Protocol.)

\item[Digest.] A SHA-256 hash over canonical or raw bytes, used as identity for comparison. Digest equality is treated as byte equality under the stated cryptographic assumption. (Chapter: Canonical Evidence Encoding.)

\item[Divergence record.] The bounded error produced when two successful workers disagree on dependency profiles: both profile hashes, at most 32 keyed difference paths, environment values and paths represented only by hashes. (Chapter: Concurrency and Failure Discipline.)

\item[Effect class.] Decision-model enum separating nonclaim research evidence from claim-bearing effect; the model's load-bearing distinction. (Chapter: Runtime Decision Model.)

\item[Exact-type identity.] The rule that identity-bearing public values compare by exact runtime type and primitive fields; equality-overriding subclasses are rejected at public interfaces. Blocks equality spoofing. (Chapters: Semantic Identity; Threat Model.)

\item[Fail-closed.] The kernel-wide posture that irregular, ambiguous, or out-of-bounds input produces an explicit failure rather than a repaired, defaulted, or partial result. (All mechanism chapters; named throughout.)

\item[Finite arrow (\texttt{FiniteArrow}).] A typed value with explicit source and target objects drawn from finite inventories; the unit of the fractal bridge's reference kernel. (Chapter: Integration Contracts, TKS section.)

\item[Fractal bridge.] The contract mapping legacy fractal digit sequences to typed arrows via a caller-supplied catalog, composing via semigroupoid paths, failing closed on missing digits or source/target mismatch; proves type-correct composition only, never digit meaning. (Chapter: Integration Contracts.)

\item[Gate 4.] The program's engineering review gate under which the v0.1 candidate and its companions were receipted. The bounded Gate-4 receipt (2026-08-04, commit a1cd9901) is the monograph's sole test evidence. (Chapter: Evidence, Obligations, and Program.)

\item[Hash-bound artifact.] An input file (source, ledger, schema, compiler bytes, fixture) whose SHA-256 is recorded in a run's manifest, closing the receipt's claim under its inputs. Fifty-four at the Gate-4 receipt. (Chapter: Evidence, Obligations, and Program.)

\item[Hermetic bundle.] The open infrastructure of obligation one: an execution environment whose interpreter, native libraries, and toolchain are content-addressed and assembled by digest. Its absence is why isolation is not a sandbox. (Chapters: Worker Protocol; Evidence, Obligations, and Program.)

\item[Hypothesis link.] The weakest link kind: a conjectured correspondence between references, recorded for investigation, carrying no reviewed standing. (Chapter: Semantic Identity.)

\item[Identifier profile.] The declared normalized/ASCII profile to which review-visible machine identifiers are confined; unrestricted prose Unicode lives in separate fields and never participates in identifier comparison. (Chapter: Canonical Evidence Encoding.)

\item[Import-DAG manifest.] The open design of obligation five: a canonical, hash-bound description of a multi-file program's full dependency graph. Until admitted, \texttt{use} imports are rejected and compatibility evidence covers single files only. (Chapters: Worker Protocol; Evidence, Obligations, and Program.)

\item[Injective rendering.] The invariant that the rendered string form of a semantic reference is injective on structural tuples: components are escaped or length-prefixed before joining, so delimiter-bearing components cannot alias distinct tuples. (Chapters: Semantic Identity; Threat Model, identifier collision.)

\item[Integrated run.] The three-suite Gate-4 execution within which the focused Core suite ran: 208 test methods, 272 parameterized subtests, 54 hash-bound artifacts. (Chapter: Evidence, Obligations, and Program.)

\item[\texttt{legacy\_only}.] The explicit marker of the sole admitted compatibility profile: verified legacy compiler execution for comparison purposes, no CAS v8 semantic translation. (Chapter: Worker Protocol.)

\item[Link kind.] The enum over cross-namespace links: \textsc{translation}, \textsc{alignment}, \textsc{hypothesis}. Deliberately excludes \textsc{identity}; no coercion path from any kind to identity exists. (Chapter: Semantic Identity.)

\item[Marked cover.] A Noetic Topos 2.0 artifact declaring an intended covering family of contexts for a target context. A declaration only: it proves no measurement true, and no operation derives a resolved statement from it. (Chapter: Integration Contracts.)

\item[Namespace.] One of the seven disjoint ontological families of the identity layer: TKS Foundation, TKS World, CAS Drive, CAS Domain, CAS Register, CAS State, Noetic Topos object. Disjointness makes cross-family conflation structurally false. (Chapter: Semantic Identity.)

\item[Namespace collision (package-path).] The distinct, repository-level failure that bounded the Gate-4 receipt: four pre-existing collisions among package paths in the wider monorepo surfaced by the repository-wide collector after 1{,}200 items, precluding any repository-wide claim. (Chapter: Evidence, Obligations, and Program.)

\item[No-silent-identity rule.] The rule that positional or historical correspondences across namespaces are typed links, never identities. The kernel's answer to the program's documented conflation incidents. (Chapters: Semantic Identity; Threat Model, namespace conflation.)

\item[Noetic Topos 2.0.] Companion candidate carrying the topos-theoretic layer; contributes the ambient-versus-resolved distinction and the marked-cover discipline. (Chapter: Integration Contracts.)

\item[Nonclaim research evidence.] The effect class that may pass a structural replay comparison when every exact compatibility field and digest validates. The only passable class in v0.1. (Chapter: Runtime Decision Model.)

\item[Open obligation.] One of the six stated gaps between the v0.1 candidate and any status change: hermetic bundle; Python/Rust conformance; v8 parser and migration; claim authority; import-DAG manifest; independent review. (Chapter: Evidence, Obligations, and Program.)

\item[Parameterized subtest.] A test outcome generated by running one assertion body over a fixture-table row; counted separately (272) from test methods (208) in the integrated run so that neither count inflates the other. (Chapter: Evidence, Obligations, and Program.)

\item[Product state.] Second field of the research boundary; \textsc{not licensed} is the only value constructible for Ring-4 material in v0.1. (Chapter: The Research Boundary.)

\item[Receipt.] Generically: a historical, digest-bound record that a stated software behavior occurred over stated inputs on a stated date. Never an experiment, a review, or authority. (Chapters: Worker Protocol; Evidence, Obligations, and Program.)

\item[Receipt laundering.] The threat of presenting structural receipts, self-hashes, or generated licenses as authority; historically instantiated by AI-authored self-hashed licensing receipts in the program's past. Blocked terminally by the constant-false invariant. (Chapter: Threat Model.)

\item[Registry.] A byte-backed ledger of records (claim registry, theorem registry) loaded only through the validating loader, with recursively immutable values and byte-bound record digests. (Chapter: Registry Integrity.)

\item[Rejection catalog.] The closed list of conditions under which the registry loader refuses input: duplicate JSON keys, unknown root or record fields, wrong enums or shapes, invalid identifiers, ambiguous versions, missing required metadata. Exhaustively fixtured. (Chapters: Registry Integrity; Reference Tables.)

\item[Replay requirement.] Protocol step eight: fresh re-execution against the same source and compiler identity is required before any structural decision can pass; stored receipts alone satisfy nothing. (Chapters: Worker Protocol; Threat Model, replay evasion.)

\item[Resolved evidence.] The aspirational cross-context form of evidence: local records glued into a statement consumable beyond their contexts. Requires restriction and gluing definitions that remain open topos-side work; nothing in the lane is resolved in v0.1. (Chapter: Integration Contracts.)

\item[Review matrix / review criteria.] The enumerated docket of what independent review must examine: 53 criteria at the Gate-4 receipt, 20 open, with author-closed criteria reopening under independence. (Chapter: Evidence, Obligations, and Program.)

\item[Ring 4.] The program's research-candidate status stratum: admission \textsc{open}, product \textsc{not licensed}, disclosure withheld. Every artifact in this monograph carries it. (Chapters: The Research Boundary; all boundary statements.)

\item[Self-digest.] A digest an object carries over its own content; detects post-construction mutation. Explicitly not a signature, reviewer identity, or authority grant. (Chapters: The Research Boundary; Threat Model, receipt laundering.)

\item[Semantic reference.] The five-field structural tuple (namespace, \texttt{local\_id}, version, \texttt{source\_id}, \texttt{definition\_sha256}) whose full-tuple structural equality is the kernel's only identity. (Chapter: Semantic Identity.)

\item[Semigroupoid path (\texttt{SemigroupoidPath}).] A nonempty finite sequence of finite arrows with source/target agreement checked at every consecutive pair; the bridge's composition witness. No identity arrows or total composition are assumed. (Chapter: Integration Contracts.)

\item[Structural consistency.] The property that an artifact's fields, digests, and boundary all validate. Never authorization: the kernel's central distinction. (Chapters: Runtime Decision Model; Threat Model.)

\item[Structural decision.] A decision-model outcome asserting replay-verified compatibility for nonclaim research evidence; the strongest verdict the v0.1 kernel can issue. (Chapter: Runtime Decision Model.)

\item[Suite test-ID hash.] The SHA-256 over the sorted, fully qualified identifiers of a suite's tests; binds a receipt to an exact test set rather than a count. The focused Core suite's is \texttt{E2AF5014DE28DD0702211EF32A2757B99B869A13FAFBB0B8D1B18A50A9F8FEF5}. (Chapter: Evidence, Obligations, and Program.)

\item[Tagged decimal.] The encoder's representation of exact decimal quantities, tagged so that $\mathrm{Decimal}(\texttt{"1.0"})$, the string \texttt{"1.0"}, and the integer $1$ encode pairwise distinctly, defeating type puns. (Chapter: Canonical Evidence Encoding.)

\item[Theorem registry.] The byte-backed ledger of formal-result records, governed by the same loader discipline as the claim registry. (Chapter: Registry Integrity.)

\item[TKS v8.0.] Companion candidate carrying TKS-side formal content; supplies and reviews the arrow catalogs consumed by the fractal bridge. (Chapter: Integration Contracts.)

\item[Translation link.] The link kind recording a reviewed cross-namespace rendering of one ontology's element in another's terms; the strongest representable cross-namespace statement, still not identity. (Chapter: Semantic Identity.)

\item[Trust root.] The authenticated basis --- authority, reviewer identities, key custody, revocation practice --- that authorization would require and that v0.1 deliberately lacks; its absence is why the kernel's answer is constant \emph{no}. (Chapters: Runtime Decision Model; Evidence, Obligations, and Program.)

\item[Type pun.] A cross-type encoding collision (e.g.\ a Decimal, string, and integer sharing bytes) that would let different typed facts become digest-identical; excluded by tagged encoding. (Chapters: Canonical Evidence Encoding; Threat Model.)

\item[Worker.] The fresh \texttt{python -I -S -B} process with sanitized environment in which verified compiler bytes execute; the unit of Python-process isolation. (Chapter: Worker Protocol.)

\item[Worker runtime binding.] The pre-execution sampling of interpreter fields, platform data, and profile-bound architecture inputs into the receipt, so the receipt records the runtime that actually executed. (Chapter: Worker Protocol.)

\end{description}
